\begin{proof}
\begin{enumerate}
\item \textbf{Notation.}  
Let
\[
\mathbf a=(a_{1},\dots ,a_{n}),\qquad 
\mathbf b=(b_{1},\dots ,b_{n})\in\mathbb R^{n},
\]
and set
\[
A=\sum_{i=1}^{n}a_{i}^{2},\qquad 
B=\sum_{i=1}^{n}b_{i}^{2},\qquad 
C=\sum_{i=1}^{n}a_{i}b_{i}.
\]

\item \textbf{Degenerate case.}  
If \(B=0\) then every \(b_{i}=0\); consequently \(C=0\) and \(AB=0\). Hence
\[
\bigl(\sum_{i=1}^{n}a_{i}b_{i}\bigr)^{2}=0=AB,
\]
so the inequality holds with equality. In the sequel we assume \(B>0\).

\item \textbf{Quadratic construction.}  
For a real parameter \(t\) consider
\[
Q(t)=\sum_{i=1}^{n}(a_{i}-t b_{i})^{2}.
\]
Each term is a square, therefore \(Q(t)\ge 0\) for all \(t\in\mathbb R\).

\item \textbf{Expansion.}  
Expanding gives
\begin{align*}
Q(t) &=\sum_{i=1}^{n}\bigl(a_{i}^{2}-2t a_{i}b_{i}+t^{2}b_{i}^{2}\bigr)\\
     &=A-2Ct+Bt^{2}.
\end{align*}
Thus \(Q(t)=Bt^{2}-2Ct+A\) is a quadratic polynomial with leading coefficient \(B>0\).

\item \textbf{Discriminant argument.}  
A quadratic \(Bt^{2}-2Ct+A\) that is non‑negative for every real \(t\) must have non‑positive discriminant:
\[
\Delta =(-2C)^{2}-4BA =4\bigl(C^{2}-AB\bigr)\le 0.
\]
Hence \(C^{2}\le AB\).

\item \textbf{Desired inequality.}  
Since \(C=\sum_{i=1}^{n}a_{i}b_{i}\), \(A=\sum_{i=1}^{n}a_{i}^{2}\) and \(B=\sum_{i=1}^{n}b_{i}^{2}\), the inequality \(C^{2}\le AB\) is exactly
\[
\Bigl(\sum_{i=1}^{n}a_{i}b_{i}\Bigr)^{2}\le
\Bigl(\sum_{i=1}^{n}a_{i}^{2}\Bigr)
\Bigl(\sum_{i=1}^{n}b_{i}^{2}\Bigr).
\]

\item \textbf{Equality condition.}  
Equality occurs iff \(\Delta =0\), i.e. \(C^{2}=AB\). With \(B>0\) this means the quadratic \(Q(t)\) has a double root, so
\[
t_{0}= \frac{C}{B}
\]
satisfies \(a_{i}=t_{0}b_{i}\) for every \(i\); that is, \(\mathbf a\) is a scalar multiple of \(\mathbf b\).  
In the degenerate case \(B=0\) we already have equality. Hence equality holds precisely when \(\mathbf b=0\) or \(\mathbf a\) and \(\mathbf b\) are linearly dependent (one is a scalar multiple of the other).

\item \textbf{Conclusion.}  
For any real sequences \((a_{i})_{i=1}^{n}\) and \((b_{i})_{i=1}^{n}\) we have
\[
\Bigl(\sum_{i=1}^{n}a_{i}b_{i}\Bigr)^{2}\le
\Bigl(\sum_{i=1}^{n}a_{i}^{2}\Bigr)
\Bigl(\sum_{i=1}^{n}b_{i}^{2}\Bigr),
\]
with equality exactly in the cases described above. \(\square\)
\end{enumerate}
\end{proof}